
\documentclass[10pt, a4paper]{article}
\usepackage[margin=0.75in]{geometry}
\usepackage{amsmath, amssymb, amsthm, mathtools}
\usepackage{enumitem}
\usepackage[T1]{fontenc}
\usepackage{lmodern}
\usepackage{microtype}
\usepackage{hyperref}
\usepackage{xcolor}
\usepackage[most,skins,breakable]{tcolorbox}
\usepackage{fancyhdr}
\usepackage{graphicx}
\usepackage{float}
\usepackage{caption}
\usepackage{subcaption}
\usepackage{tikz}
\usepackage{pgfplots}
\pgfplotsset{compat=1.18}
\usetikzlibrary{arrows.meta, positioning, calc, shapes.geometric,
  decorations.pathmorphing, patterns, automata, trees, fit, backgrounds, matrix}

\hypersetup{colorlinks=true, linkcolor=blue!60!black}
\setlength{\parindent}{0pt}
\setlength{\parskip}{3pt}
\setlist{nosep, leftmargin=1.5em}

\pagestyle{fancy}
\fancyhf{}
\lhead{{\small MATH 999 \quad Chapter 1 Practice Problems}}
\rhead{{\small Fall 2099}}
\cfoot{\thepage}

\definecolor{easygreen}{HTML}{1E8449}
\definecolor{easygreenbg}{HTML}{D5F5E3}
\definecolor{medorange}{HTML}{B9770E}
\definecolor{medorangebg}{HTML}{FDEBD0}
\definecolor{hardred}{HTML}{922B21}
\definecolor{hardredbg}{HTML}{FADBD8}
\definecolor{solgray}{HTML}{2C3E50}
\definecolor{solgraybg}{HTML}{EAECEE}

\newtcolorbox{qeasy}[1][]{%
  enhanced, breakable,
  colback=easygreenbg, colframe=easygreen, coltitle=white,
  fonttitle=\bfseries, title={#1},
  arc=1.5mm, boxrule=0.8pt,
  left=4pt, right=4pt, top=3pt, bottom=3pt, after skip=8pt,
  drop shadow southeast}
\newtcolorbox{qmedium}[1][]{%
  enhanced, breakable,
  colback=medorangebg, colframe=medorange, coltitle=white,
  fonttitle=\bfseries, title={#1},
  arc=1.5mm, boxrule=0.8pt,
  left=4pt, right=4pt, top=3pt, bottom=3pt, after skip=8pt,
  drop shadow southeast}
\newtcolorbox{qhard}[1][]{%
  enhanced, breakable,
  colback=hardredbg, colframe=hardred, coltitle=white,
  fonttitle=\bfseries, title={#1},
  arc=1.5mm, boxrule=0.8pt,
  left=4pt, right=4pt, top=3pt, bottom=3pt, after skip=8pt,
  drop shadow southeast}
\newtcolorbox{solution}[1][]{%
  enhanced, breakable,
  colback=solgraybg, colframe=solgray, coltitle=white,
  fonttitle=\bfseries, title={#1},
  arc=1.5mm, boxrule=0.8pt,
  left=4pt, right=4pt, top=3pt, bottom=3pt, after skip=8pt,
  drop shadow southeast}

\newcommand{\R}{\mathbb{R}}
\newcommand{\N}{\mathbb{N}}
\newcommand{\Q}{\mathbb{Q}}
\newcommand{\Z}{\mathbb{Z}}
\newcommand{\C}{\mathbb{C}}
\newcommand{\eps}{\varepsilon}

\title{\textbf{MATH 999: Introduction to Widget Theory}\\[4pt]
  {\Large Chapter 1 Practice Problems}}
\author{{\small Exemplar University \quad Fall 2099}}
\date{}

\begin{document}
\maketitle
\newpage

% ═══════════════════════════════════════════════════════════════════════
% PROBLEMS
% ═══════════════════════════════════════════════════════════════════════

\section*{Problems}

\begin{qeasy}[Q1]
Compute $\|w\|_W$ for $w = (2, 0, 6, 0, 0, \ldots)$.
\end{qeasy}

\begin{qeasy}[Q2]
Let $w = (1, 1, 1, \ldots)$. Compute $\|w\|_W$.
\end{qeasy}

\begin{qeasy}[Q3]
True or false: $\|w\|_W = 0$ implies $w = (0, 0, 0, \ldots)$.
Justify.
\end{qeasy}

\begin{qmedium}[Q4]
Let $w^{(n)}_k = n/(n+k)$. Show $\|w^{(n)}\|_W \le 1$ for
all $n$ and find $\lim_{n \to \infty} \|w^{(n)}\|_W$.
\end{qmedium}

\begin{qmedium}[Q5]
Prove $\|u + v\|_W \le \|u\|_W + \|v\|_W$.
\end{qmedium}

\begin{qmedium}[Q6]
State the Widget Approximation Theorem. What are its hypotheses
and conclusions?
\end{qmedium}

\begin{qmedium}[Q7]
Give an example of a bounded widget sequence that does not
converge but has a convergent subsequence.
\end{qmedium}

\begin{qhard}[Q8]
Prove the Tail Bound Lemma: if $\|w\|_W \le M$, then
$\sup_{k > K} |w_k| \cdot 2^{-k} \le M$ for all $K \ge 1$.
\end{qhard}

\begin{qhard}[Q9]
Prove: $\|w^{(n)} - w\|_W \to 0$ implies $w^{(n)}_k \to w_k$
for every $k$. Is the converse true?
\end{qhard}

\begin{qhard}[Q10]
Prove the Widget Approximation Theorem.
\end{qhard}


\newpage

% ═══════════════════════════════════════════════════════════════════════
% SOLUTIONS
% ═══════════════════════════════════════════════════════════════════════

\section*{Solutions}

\begin{solution}[Q1]
$w = (2, 0, 6, 0, 0, \ldots)$.
\begin{align*}
  |w_1| \cdot 2^{-1} &= 1, \\
  |w_3| \cdot 2^{-3} &= 6/8 = 3/4.
\end{align*}
All other terms are 0. So $\|w\|_W = \boxed{1}$.
\end{solution}

\begin{solution}[Q2]
$|w_k| \cdot 2^{-k} = 2^{-k}$, maximized at $k = 1$.

$\|w\|_W = 2^{-1} = \boxed{1/2}$.
\end{solution}

\begin{solution}[Q3]
True. $\|w\|_W = 0$ means $|w_k| \cdot 2^{-k} \le 0$ for all $k$.
Since $2^{-k} > 0$, this forces $w_k = 0$ for all $k$.
\end{solution}

\begin{solution}[Q4]
$\frac{n}{n+k} \cdot 2^{-k} \le 2^{-k} \le 1/2$ for all $k \ge 1$.
So $\|w^{(n)}\|_W \le 1$.

At $k = 1$: $\frac{n}{n+1} \cdot \frac{1}{2} \to \frac{1}{2}$.
For $k \ge 2$: the terms $\to 2^{-k} \le 1/4 < 1/2$.

$\lim \|w^{(n)}\|_W = \boxed{1/2}$.
\end{solution}

\begin{solution}[Q5]
For each $k$:
\[
  |u_k + v_k| \cdot 2^{-k} \le (|u_k| + |v_k|) \cdot 2^{-k}
  \le \|u\|_W + \|v\|_W.
\]
Taking the sup over $k$ gives $\|u + v\|_W \le \|u\|_W + \|v\|_W$. $\qed$
\end{solution}

\begin{solution}[Q6]
\textit{Statement:} Let $(w^{(n)})$ satisfy $\|w^{(n)}\|_W \le M$ for all $n$.
Then there exists a subsequence $(w^{(n_j)})$ and $w^* \in W$ with
$\|w^{(n_j)} - w^*\|_W \to 0$.

Hypotheses: the sequence is bounded in the widget norm.
Conclusion: a norm-convergent subsequence exists.
\end{solution}

\begin{solution}[Q7]
$w^{(n)}_k = (-1)^n$ if $k = 1$, and $w^{(n)}_k = 0$ for $k \ge 2$.

Then $\|w^{(n)}\|_W = 1/2$ for all $n$ (bounded), but $w^{(n)}_1 = (-1)^n$
does not converge. Taking the subsequence $n = 2j$ gives $w^{(2j)} \to w^*$
with $w^*_1 = 1$, $w^*_k = 0$ for $k \ge 2$.
\end{solution}

\begin{solution}[Q8]
$\sup_{k > K} |w_k| \cdot 2^{-k}$ is a supremum over a subset of
$\{|w_k| \cdot 2^{-k} : k \ge 1\}$, so it is at most the full sup, which
is $\|w\|_W \le M$. $\qed$
\end{solution}

\begin{solution}[Q9]
$(\Rightarrow)$ Fix $k$. Then $|w^{(n)}_k - w_k| \cdot 2^{-k} \le
\|w^{(n)} - w\|_W \to 0$. Since $2^{-k} > 0$:
$|w^{(n)}_k - w_k| \le 2^k \|w^{(n)} - w\|_W \to 0$.

Converse is false. Let $w^{(n)}_k = 2^k$ if $k = n$, and $0$ otherwise.
Then $w^{(n)}_k \to 0$ for each fixed $k$. But
$\|w^{(n)}\|_W = 2^n \cdot 2^{-n} = 1 \not\to 0$.
\end{solution}

\begin{solution}[Q10]
We construct $w^*$ component by component.

\textit{Step 1.} $(w^{(n)}_1)$ is bounded by $2M$. By Bolzano-Weierstrass,
extract a subsequence $(w^{(n^1_j)})$ with $w^{(n^1_j)}_1 \to w^*_1$.

\textit{Step 2.} From $(w^{(n^1_j)})$, extract a further subsequence
$(w^{(n^2_j)})$ with $w^{(n^2_j)}_2 \to w^*_2$.

\textit{Step 3.} Continue for all $k$. The diagonal subsequence
$w^{(n^j_j)}$ converges in every component.

\textit{Step 4.} Let $\eps > 0$. Choose $K$ with $M \cdot 2^{-K} < \eps/2$.
For $j$ large enough:
\[
  \|w^{(n^j_j)} - w^*\|_W
  \le \max\!\left(\max_{k \le K} |w^{(n^j_j)}_k - w^*_k| \cdot 2^{-k},\;
       M \cdot 2^{-K}\right) < \eps. \qed
\]
\end{solution}


\end{document}
